% Part: incompleteness
% Chapter: introduction
% Section: decidability

\documentclass[../../../include/open-logic-section]{subfiles}

\begin{document}

\olfileid[tr]{inc}{int}{dec}

\olsection{Karar Verilemezlik ve Eksiklik}

G\"odel'in eksiklik teoremleri ispatı sözdizimin aritmetikleştirilmesini
gerektirir. Fakat bu olmadan da, yalnızca bir kuramın belirli bağıntıları
!!{represents} ettiği ve bunların tam olarak !!{decidable} bağıntılar
olduğu varsayımıyla bazı güzel sonuçlar elde
edebiliriz. İspat, durma probleminin karar verilemezliği ispatına benzer bir
köşegen savıdır.

\begin{thm}
$\Gamma$ her bağıntıyı !!{represents} eden tutarlı bir kuramsa---burada
``her bağıntı'' ile her !!{decidable} bağıntı kastedilmektedir---$\Gamma$
!!{decidable} olma özelliğini taşımaz.
\end{thm}

\begin{proof}
$\Gamma$'nın !!{decidable} olduğunu varsayalım. $\Gamma$ her bağıntıyı
!!{represents} ediyorsa---yine !!{decidable} bağıntıları kastediyoruz---
tutarsız olmak zorunda olduğunu göstereceğiz.

!!^{decidable} özellikler (tek yerli bağıntılar), tek bir serbest değişkeni
olan !!{formula}s{}le temsil edilir. $!A_0(x)$, $!A_1(x)$, \dots, böyle
bütün !!{formula}s{}in hesaplanabilir bir sayımı olsun. Şimdi aşağıdaki
$D \subseteq \Nat$ kümesini ele alın:
\[
D = \Setabs{n}{\Gamma \Proves \lnot !A_n(\num{n})}
\]
$D$ kümesi !!{decidable} olur; çünkü $n \in D$ olup olmadığını önce
$!A_n(x)$'i ve bundan da $\lnot !A_n(\num{n})$'i hesaplayarak sınayabiliriz.
% [tr source emendation] Upstream writes !A rather than the indexed !A_n
% in the next two occurrences; the defined diagonal family requires !A_n.
Açıkça $!A_n(x)$ içindeki her serbest $x$ geçişi yerine $\num{n}$ terimini
koymak ve $!A_n(\num{n})$'in başına $\lnot$ getirmek mekanik bir iştir.
Varsayım gereği $\Gamma$ !!{decidable} olduğundan
$\lnot !A_n(\num{n}) \in \Gamma$ olup olmadığını sınayabiliriz. Öyleyse
$n \in D$; değilse $n \notin D$ olur. Dolayısıyla $D$ de !!{decidable} olur.

$\Gamma$ bütün özellikleri !!{represents} ettiğine ve söz konusu özellikler
!!{decidable} olduğuna göre $D$'yi de !!{represents} eder. $D$'yi
$\Gamma$ içinde temsil eden !!{formula}s da $!A_0(x)$, $!A_1(x)$, \dots
arasında yer alır ve bu kümeyi !!{represents} eder. Öyleyse $!A_d(x)$'in
$D$'yi $\Gamma$ içinde !!{represents} edeceği bir $d$ sayısı
seçin. $d \notin D$ ise $!A_d(x)$, $D$'yi !!{represents} ettiğinden
$\Gamma \Proves \lnot !A_d(\num{d})$ olur. Fakat bu, $d$'nin $D$'nin
tanımlayıcı koşulunu karşıladığı ve dolayısıyla $d \in D$ olduğu anlamına
gelir. Bu, $d \notin D$ ile çelişir. Dolaylı ispatla $d \in D$'dir.

$d \in D$ olduğundan $D$'nin tanımı gereği
$\Gamma \Proves \lnot !A_d(\num{d})$ olur. Öte yandan $!A_d(x)$, $D$'yi
$\Gamma$ içinde !!{represents} ettiğinden $\Gamma \Proves !A_d(\num{d})$
olur. O hâlde $\Gamma$ tutarsızdır.
\end{proof}

\begin{explain}
Önceki teorem, bütün bağıntıları !!{represents} eden hiçbir tutarlı kuramın,
bu bağıntılar !!{decidable} olduğunda, !!{decidable} olamayacağını gösterir.
$\Th{Q}$'nun bütün bağıntıları !!{represents} ettiğini---burada
!!{decidable} bağıntıları kastediyoruz---göstereceğiz; bu,
$\Th{PA}$ ve $\Th{TA}$ gibi $\Th{Q}$'yu içeren bütün kuramların da bunları
temsil ettiği ve dolayısıyla !!{decidable} olmadığı anlamına gelir.
(Bu kuramların hepsi standart modelde doğru olduğundan hepsi tutarlıdır.)

Bu sonucu birinci eksiklik teoreminin zayıf bir sürümünü elde etmek için de
kullanabiliriz. !!{axiomatizable} ve !!{complete} olan her kuram
!!{decidable} olur. Öyleyse !!{axiomatizable} olan ve bütün özellikleri
!!{represents} eden (bu özellikler !!{decidable} olur) tutarlı kuramlar
!!{complete} olamaz.
\end{explain}

\begin{thm}
$\Gamma$ !!{axiomatizable} ve !!{complete} ise !!{decidable} olur.
\end{thm}

\begin{proof}
Her tutarsız kuram !!{decidable} olur; çünkü tutarsız kuramlar bütün
!!{sentence}s{}i içerir. Dolayısıyla ``$!A \in \Gamma$ mı?'' sorusunun
yanıtı her zaman ``evet''tir ve böylece karar verilebilir.

Şimdi $\Gamma$'nın tutarlı, ayrıca !!{axiomatizable} ve !!{complete}
olduğunu varsayalım. $\Gamma$ !!{axiomatizable} olduğundan
!!{computably enumerable} olur. Çünkü $\Gamma$'nın aksiyomlarından bütün
doğru !!{derivation}s{}i hesaplanabilir bir fonksiyonla sayabiliriz. Doğru
bir !!{derivation}{}den onun ürettiği !!{sentence}{}yi hesaplayabiliriz; bu
tümceyi o türetim !!{derive} eder;
böylece bunlar birlikte $\Gamma$'nın bütün teoremlerini sayan hesaplanabilir
bir fonksiyon verir. Bir !!{sentence}, ancak ve ancak $\lnot !A$ bir teorem
değilse $\Gamma$'nın bir teoremidir; çünkü $\Gamma$ tutarlı ve
!!{complete}{}dır. Dolayısıyla $!A \in \Gamma$ olup olmadığına şöyle karar
verebiliriz. $\Gamma$'nın bütün teoremlerini sayın. Bu listede $!A$
belirdiğinde $\Gamma \Proves !A$ olduğunu biliriz. Listede $\lnot !A$
belirdiğinde $\Gamma \Proves/ !A$ olduğunu biliriz. $\Gamma$ !!{complete}
olduğundan bu durumlardan biri sonunda gerçekleşir; böylece yordam sonunda
bir yanıt üretir.
\end{proof}

\begin{cor}
\ollabel{cor:incompleteness}
$\Gamma$ tutarlı ve !!{axiomatizable} olup her özelliği !!{represents}
ediyorsa---bu özellikler !!{decidable} sayılır---!!{complete} değildir.
\end{cor}

\begin{proof}
$\Gamma$ !!{complete} olsaydı önceki teorem uyarınca !!{decidable} olurdu;
çünkü tutarlı ve !!{axiomatizable} bir kuramdır. Fakat $\Gamma$ her özelliği
!!{represents} ettiğinden ve bu özellikler !!{decidable} olduğundan, ilk
teorem uyarınca !!{decidable} olma özelliğini taşımaz.
\end{proof}

\begin{prob}
$\Th{TA} = \Setabs{!A}{\Sat{N}{!A}}$'nın !!{axiomatizable} olmadığını
gösterin. $\Th{TA}$'nın bütün karar verilebilir özellikleri temsil ettiğini
varsayabilirsiniz.
\end{prob}

Örneğin $\Th{Q}$'nun bütün özellikleri !!{represents} ettiğini---burada
!!{decidable} özellikleri kastediyoruz---kurduğumuzda, sonuç $\Th{Q}$'nun
eksik olmak zorunda olduğunu söyler. Ancak
ispatı bağımsız bir !!{sentence} örneği vermez; yalnızca böyle !!a{sentence}{}nin
var olmak zorunda olduğunu gösterir. Bunun için sözdizimi aritmetikleştirmeli
ve G\"odel'in özgün ispat düşüncesini izlemeliyiz. Elbette ilk iddiayı, yani
$\Th{Q}$'nun gerçekten bütün özellikleri !!{represents} ettiğini---burada da
!!{decidable} özellikleri kastediyoruz---hâlâ göstermemiz gerekir.

Her \emph{ilginç} kuramın eksik veya karar verilemez olmadığı belirtilmelidir.
İlginç matematiksel olguları betimleyecek kadar güçlü olup G\"odel sonucunun
koşullarını sağlamayan birçok kuram vardır. Örneğin standart modelde doğru
olan, aritmetik dilinin $\times$ içermeyen !!{sentence}s{}i kümesi
$\Th{Pres} = \Setabs{!A \in \Lang{L_{A^+}}}{\Sat{N}{!A}}$ hem tam hem de
karar verilebilirdir. Bu kurama Presburger aritmetiği denir; yalnızca
$\Obj{0}$, $\prime$ ve $+$ ile kurulabilen doğal sayılar hakkındaki bütün
doğruları ispatlar.

\end{document}
